\documentclass[../../main.tex]{subfiles}
\begin{document}

\begin{flushright}
  \footnotesize\textit{【自動文字起こし・要確認】}
\end{flushright}

{\bf [解]} (A)から，$g(x)=f(x)-x$ とすると，$g(1)=g(-1)=0$ かつ $g$は2次関数だから，
\begin{align}
  g(x) = a(x+1)(x-1)
\end{align}
とおけるので，$f(x) = a(x+1)(x-1)+x$ である．したがってBから，
\begin{align}
  f(x) &\le 3x^2-1 \\
  a(x^2-1) &\le 3x^2-x-1 \cdots (*)
\end{align}
$x=\pm 1$ の時は明らかに成立するので，以下 $|x| < 1$ とする．$x^2-1 < 0$ から，
\begin{align}
  a \ge \dfrac{3x^2-x-1}{x^2-1} \equiv h(x) \cdots (1)
\end{align}
となる．$h(x) = 3 + \dfrac{-x+2}{x^2-1}$ だから，
\begin{align}
  h'(x) = \dfrac{-(x^2-1) - 2x(-x+2)}{(x^2-1)^2} = \dfrac{x^2-4x+1}{(x^2-1)^2}
\end{align}
となり，下表をうる．
\begin{table}[htb]
  \centering
  \begin{tabular}{c|ccccc}
    $x$ & $-1$ & $\cdots$ & $2-\sqrt{3}$ & $\cdots$ & $1$ \\ \hline
    $h'$ & & $+$ & $0$ & $-$ & \\ \hline
    $h$ & $-\infty$ & $\nearrow$ & $2-\dfrac{\sqrt{3}}{2}$ & $\searrow$ & $-\infty$
  \end{tabular}
  \caption{$h(x)$ の増減表}
\end{table}
したがって，(1)が任意の $|x|<1$ で成立する $a$ の条件は
\begin{align}
  a \ge 2-\dfrac{\sqrt{3}}{2} \cdots (2)
\end{align}
である．このもとで $I$ の $\min$ をもとめれば良い．$f'(x) = 2ax+1$ だから，
\begin{align}
  I &= \int_{-1}^1 (2ax+1)^2 dx = 2 \int_0^1 (4a^2 x^2 + 1) dx \\
  &= 2 \left[ \dfrac{4}{3}a^2 x^3 + x \right]_0^1 = 2 \left( \dfrac{4}{3}a^2 + 1 \right) \cdots (3)
\end{align}
(2)から $a^2 \ge \dfrac{19-8\sqrt{3}}{4}$ だから，(3)より
\begin{align}
  I \ge \dfrac{4}{3}(11-4\sqrt{3})
\end{align}

\bigskip
\noindent{\bf [解2]} (以下略)
$h(x) = (3-a)x^2 - x + a - 1$ とおく．$a=3$ の時 $h(x) = -x + 2$ となり，$-1 \le x \le 1$ では $h(x) \ge 0$ だから不成立．以下 $f(1) = 1, f(-1) = 3$ に注意する．$h(x) = 0$ の判別式を $D$ とする．

$1^\circ \ 3 > a$ の時
\begin{align}
  D \le 0 \quad \text{or} \quad \left\{ \dfrac{1}{2(3-a)} \le -1, \ 1 \le \dfrac{1}{2(3-a)} \right\} \iff \dfrac{4-\sqrt{3}}{2} \le a < 3
\end{align}
$2^\circ \ 3 < a$ の時
$f(1) > 0, f(-1) > 0$ から成立

以上から $\dfrac{4-\sqrt{3}}{2} \le a$ である（以下略）

\end{document}
